\section{The Intensional Turn: From Extension to Conceptualization}

The fundamental inadequacy of extensional structures for capturing conceptual content emerges clearly in Guarino's analysis. A conceptualization must remain invariant across different world states, yet extensional relations vary precisely with these states. This tension motivates the crucial turn toward intensionality.

\begin{definition}[World and World State]
	With respect to a specific system $S$ we want to model, a \emph{world state} for $S$ is a maximal observable state of affairs, i.e., a unique assignment of values to all the observable variables that characterize the system. A \emph{world} is a totally ordered set of world states, corresponding to the system's evolution in time. If we abstract from time for the sake of simplicity, a world state coincides with a world.
\end{definition}

This definition embeds a critical methodological choice: we fix a granularity of observation and declare certain variable configurations ``maximal.'' The epistemic humility here is noteworthy—we do not claim to capture all possible states of $S$, but rather all states \emph{observable at our chosen level of analysis}.

\begin{definition}[Intensional Relation or Conceptual Relation]
	Let $S$ be an arbitrary system, $D$ an arbitrary set of distinguished elements of $S$, and $W$ the set of world states for $S$ (also called worlds, or possible worlds). The pair $\langle D, W \rangle$ is called a \emph{domain space} for $S$, as it intuitively fixes the space of variability of the universe of discourse $D$ with respect to the possible states of $S$. An \emph{intensional relation} (or \emph{conceptual relation}) $\rho^n$ of arity $n$ on $\langle D, W \rangle$ is a total function
	$$\rho^n : W \to 2^{D^n}$$
	from the set $W$ into the set of all $n$-ary (extensional) relations on $D$.
\end{definition}

\subsection{The Carnapian Heritage and Modal Semantics}

This definition instantiates the classical Carnapian approach to intension: an intension is a function from possible worlds to extensions \cite{Carnap1956}. The conceptual relation $\rho^n$ does not \emph{vary} with worlds; rather, it is the very function that \emph{determines} how extension varies with worlds. This is the formal correlate of conceptual stability across counterfactual variation.

The significance becomes apparent when we contrast this with Backlund's observation that different definitional approaches to ``system'' often disagree about whether certain properties are essential or accidental \cite{Backlund2000a}. A purely extensional definition cannot adjudicate such disputes, for it has no resources to express modal distinctions. The intensional framework, by contrast, builds modality into its foundational architecture.

\begin{example}[Cooperates-With as Conceptual Relation]
	Following Guarino's human resources example, consider the conceptual relation \textit{cooperates-with}$^2$. Let:
	\begin{align*}
		D &= \{\text{employee}_1, \text{employee}_2, \ldots, \text{employee}_{50000}\} \\
		W &= \{\text{possible organizational states}\}
	\end{align*}
	
	The conceptual relation is not merely the set of pairs currently cooperating, but rather the function that specifies, for each possible organizational configuration $w \in W$, which pairs would be cooperating in that configuration. 
	
	Crucially, this function encodes our \emph{concept} of cooperation—what counts as cooperation—which remains stable even as the actual instantiations vary. In world $w_1$, perhaps cooperation requires merely shared goals. In this world:
	$$\text{cooperates-with}^2(w_1) = \{(\text{emp}_i, \text{emp}_j) : \text{emp}_i \text{ and } \text{emp}_j \text{ share goals in } w_1\}$$
	
	But we might imagine an alternative concept where cooperation requires both shared goals \emph{and} coordinated action. This would be a \emph{different conceptual relation} $\rho'^2$, even though it might happen to have the same extension in certain worlds.
\end{example}

\begin{definition}[Intensional Relational Structure or Conceptualization]
	An \emph{intensional relational structure} (or a \emph{conceptualization} according to Guarino) is a triple $\mathcal{C} = (D, W, \mathfrak{R})$ with:
	\begin{itemize}
		\item $D$ a set (universe of discourse)
		\item $W$ a set (possible worlds or world states)
		\item $\mathfrak{R}$ a set of conceptual relations on the domain space $\langle D, W \rangle$
	\end{itemize}
\end{definition}

\subsection{The Invariance Problem and Systems Theory}

This definition crystallizes what we might call the \emph{invariance requirement}: a conceptualization specifies the intensions—the conceptual relations—that remain fixed while extensions vary with world states. This directly addresses von Bertalanffy's concern with identifying genuine systemic properties that transcend particular states.

When von Bertalanffy characterizes a system through differential equations:
$$\frac{dQ_i}{dt} = f_i(Q_1, Q_2, \ldots, Q_n)$$
he is implicitly specifying conceptual relations—the functional dependencies $f_i$ that determine how changes in any variable $Q_i$ depend on all variables. These functions do not vary with the current values of the $Q_i$; rather, they determine how those values evolve. In Guarino's framework, each $f_i$ corresponds to a conceptual relation specifying, for each world state (configuration of $Q$ values), what the derivative relationships are.

\begin{example}[Von Bertalanffy's Biological Systems]
	Consider a simple metabolic system with states characterized by substrate concentration $S$, enzyme concentration $E$, and product concentration $P$. The universe of discourse is:
	$$D = \{\text{substrate molecules, enzyme molecules, product molecules}\}$$
	
	World states are assignments of concentrations:
	$$w \in W \text{ assigns values to } (S(w), E(w), P(w))$$
	
	A conceptual relation might be the Michaelis-Menten kinetics:
	$$\text{reaction-rate}^3 : W \to 2^{D^3}$$
	where for each world state $w$, this relation specifies which triples of (substrate, enzyme, product) entities satisfy the kinetic equation at that state. The relation itself—the functional form of Michaelis-Menten kinetics—is the invariant concept, while its extension varies with concentration states.
\end{example}